% !TEX root = ../main.tex
%% Corrigés — Chapitre 13 — Orthogonalité et distances
%% Tous les résultats ont été vérifiés par programme (calcul exact).

\corrigesChapitre{13}{Orthogonalité et distances}

\begin{deuxcolonnes}

\begin{corrige}{1}{Produit scalaire en repère orthonormé}
\begin{questions}
  \item $\vec u\cdot\vec v=2-4+6=\res{4}$ ; $\vec u\cdot\vec w=-6+0+6=\res{0}$ : $\vec u$ et $\vec w$ sont
    $\res{\text{orthogonaux}}$.
  \item $\norme{\vec u}=\res{\sqrt{14}}$, $\norme{\vec v}=\res{\sqrt{21}}$ et\\
    $\cos\theta=\frac{4}{\sqrt{14}\sqrt{21}}=\frac{2\sqrt6}{21}\approx0{,}233$, soit $\theta\approx\res{77^\circ}$.
  \item $\vec u+\vec v\,(3\pv3\pv5)$ : $9+9+25=43$, et $14+2\times4+21=43$.
  \item $\vect{AB}\,(2\pv1\pv-2)$, donc $AB=\sqrt{4+1+4}=\res{3}$.
\end{questions}
\end{corrige}

\begin{corrige}{2}{Produit scalaire dans un cube}
\begin{questions}
  \item $\vect{AC}\cdot\vect{AG}=\vect{AC}\cdot\vect{AC}+\vect{AC}\cdot\vect{CG}$. Or $(CG)$ est orthogonale au plan $(ABC)$,
    donc $\vect{AC}\cdot\vect{CG}=0$, et $AC^2=2$ : $\res{\vect{AC}\cdot\vect{AG}=2}$.
  \item $\vect{AC}\,(1\pv1\pv0)$, $\vect{AG}\,(1\pv1\pv1)$ : $1+1+0=2$.
  \item $\vect{AB}\cdot\vect{AG}=\res{1}$ ; $\vect{EC}\,(1\pv1\pv-1)$ et $\vect{BH}\,(-1\pv1\pv1)$ : $\vect{EC}\cdot\vect{BH}=\res{-1}$.
  \item $\cos\widehat{CAG}=\frac{2}{\sqrt2\times\sqrt3}=\frac{\sqrt6}{3}\approx0{,}816$, soit $\widehat{CAG}\approx\res{35^\circ}$.
\end{questions}
\end{corrige}

\begin{corrige}{3}{Vecteur normal et équation d'un plan}
\begin{questions}
  \item $\vect{AB}\,(-2\pv2\pv3)$ et $\vect{AC}\,(-3\pv-1\pv2)$ ne sont pas colinéaires (rapports $\frac32$ et $-\frac12$) : les points
    définissent un plan.
  \item $\vec n\cdot\vect{AB}=-14-10+24=0$ et $\vec n\cdot\vect{AC}=-21+5+16=0$ : $\vec n$ est orthogonal à deux vecteurs non colinéaires
    du plan, il lui est $\res{\text{normal}}$.
  \item $7x-5y+8z+d=0$ et $A$ donne $14-5-8+d=0$ :\\ $\res{7x-5y+8z-1=0}$. Pour $B$ : $-15+16-1=0$ ; pour $C$ : $-7+8-1=0$.
\end{questions}
\end{corrige}

\begin{corrige}{4}{Produits scalaires dans un tétraèdre régulier}
\begin{questions}
  \item $\vect{AB}\cdot\vect{AC}=a\times a\times\cos60^\circ=\res{\tfrac{a^2}{2}}$.
  \item $\vect{BI}=-\frac12\vect{AB}$, $\vect{AJ}=\frac12\vect{AC}$ : $\vect{BI}\cdot\vect{AJ}=-\frac14\times\frac{a^2}2=\res{-\tfrac{a^2}{8}}$.
    $\vect{IK}=\frac12\vect{BD}$ et $\vect{JK}=\frac12\vect{CD}$ (théorème des milieux) :\\
    $\vect{IK}\cdot\vect{JK}=\frac14\vect{DB}\cdot\vect{DC}$\\ $\phantom{\vect{IK}\cdot\vect{JK}}=\res{\tfrac{a^2}{8}}$.
  \item $\vect{BC}\cdot\vect{AD}=\vect{AC}\cdot\vect{AD}-\vect{AB}\cdot\vect{AD}$, soit $\frac{a^2}2-\frac{a^2}2=0$, et
    $\vect{IJ}=\frac12\vect{BC}$ donne $\res{\vect{IJ}\cdot\vect{AD}=0}$.
  \item Deux arêtes opposées sont $\res{\text{orthogonales}}$.
\end{questions}
\end{corrige}

\begin{corrige}{5}{Calculer un angle dans un cube}
\begin{questions}
  \item $H(0\pv1\pv1)$, $M\left(0\pv\frac12\pv0\right)$, $N\left(\frac12\pv0\pv1\right)$.
  \item $\vect{MH}\left(0\pv\frac12\pv1\right)$, $\vect{MN}\left(\frac12\pv-\frac12\pv1\right)$ :\\
    $\vect{MH}\cdot\vect{MN}=-\frac14+1=\res{\tfrac34}$,\\ $MH=\res{\tfrac{\sqrt5}{2}}$ et $MN=\res{\tfrac{\sqrt6}{2}}$.
  \item $\cos\widehat{HMN}=\frac{3/4}{\sqrt{30}/4}=\frac{3}{\sqrt{30}}\approx0{,}548$, d'où $\widehat{HMN}\approx\res{57^\circ}$.
\end{questions}
\end{corrige}

\begin{corrige}{6}{Équations cartésiennes de plans}
\begin{questions}
  \item $x-2y+2z+d=0$ et $A$ : $2+2+d=0$, soit $\res{x-2y+2z-4=0}$.
  \item Avec $I$ milieu de $[AB]$ :\\ $MA^2-MB^2$\\ $=\big(\vect{MA}-\vect{MB}\big)\cdot\big(\vect{MA}+\vect{MB}\big)$\\
    $=2\,\vect{BA}\cdot\vect{MI}$.\\ Donc $MA=MB\iff\vect{IM}\cdot\vect{AB}=0$. Ici $I(2\pv1\pv1)$ et $\vect{AB}\,(2\pv-2\pv-4)$ :
    $\res{x-y-2z+1=0}$.
  \item $\vec n\,(2\pv1\pv-1)$ ; $2-2-3+3=0$ : $M\in\mathcal{P}$. Les vecteurs $(1\pv-2\pv0)$ et $(0\pv1\pv1)$ sont orthogonaux à
    $\vec n$ et non colinéaires : ils dirigent $\mathcal{P}$.
  \item $\vec n\cdot(1\pv-1\pv1)=0$ : $\mathcal{P}\perp\mathcal{Q}$. Le vecteur $(-4\pv-2\pv2)=-2\vec n$ : $\mathcal{P}$ et
    $\mathcal{R}$ sont parallèles, et même $\res{\text{strictement}}$ car $M\notin\mathcal{R}$ ($-4+4+6+1\neq0$).
\end{questions}
\end{corrige}

\begin{corrige}{7}{Distance d'un point à un plan}
\begin{questions}
  \item $6-1+8-1=12\neq0$ : $M\notin\mathcal{P}$.
  \item Cette droite est dirigée par $\vec n\,(2\pv-1\pv2)$ : ses points sont les $M+t\vec n$.
  \item $2(3+2t)-(1-t)+2(4+2t)-1$\\ $=12+9t=0$, d'où $t=-\frac43$ et\\ $\res{H\left(\tfrac13\pv\tfrac73\pv\tfrac43\right)}$.
    Puis $MH=\frac43\norme{\vec n}=\frac43\times3=\res{4}$.
  \item $0+1+0-1=0$ : $Q\in\mathcal{P}$, et $MQ=\sqrt{9+4+16}=\sqrt{29}\approx5{,}4>4$ : c'est cohérent, $MH$ est la plus petite
    distance.
\end{questions}
\end{corrige}

\begin{corrige}{8}{Projeté orthogonal sur une droite}
\begin{questions}
  \item $\vect{MH_t}\,(t-3\pv2t\pv2t-3)$.
  \item $(t-3)+4t+2(2t-3)=9t-9=0$ : $t=1$ et $\res{H(1\pv3\pv3)}$.
  \item $\vect{MH}\,(-2\pv2\pv-1)$ : $\res{MH=3}$.
  \item $(t-3)^2+4t^2+(2t-3)^2=9t^2-18t+18$\\ $=9(t-1)^2+9\geq9$, avec égalité pour $t=1$ seulement : la distance est minimale en $H$.
\end{questions}
\end{corrige}

\end{deuxcolonnes}

%% ---------------------------------------------------------------------
\corrigesSouples
\begin{corrige}{9}{Un angle variable dans un cube}
\begin{questions}
  \item $B(1\pv1\pv0)$, $E(1\pv0\pv1)$, $F(1\pv1\pv1)$, et $\vect{DM}=x\vect{DF}$ donne $\res{M(x\pv x\pv x)}$.
  \item $\vect{MB}\,(1-x\pv1-x\pv-x)$ et $\vect{ME}\,(1-x\pv-x\pv1-x)$ : $MB^2=ME^2=2(1-x)^2+x^2=3x^2-4x+2$, donc $MB=ME$ (le
    triangle $MEB$ est isocèle en $M$). Puis $\vect{MB}\cdot\vect{ME}=(1-x)^2-2x(1-x)=(1-x)(1-3x)=3x^2-4x+1$, et
    $\cos\theta(x)=\frac{\vect{MB}\cdot\vect{ME}}{MB\times ME}=\frac{3x^2-4x+1}{3x^2-4x+2}$.
  \item Le discriminant de $3x^2-4x+2$ vaut $-8<0$ : ce trinôme est toujours strictement positif. En écrivant
    $f(x)=1-\frac{1}{3x^2-4x+2}$, on obtient $f'(x)=\frac{6x-4}{(3x^2-4x+2)^2}=\frac{2(3x-2)}{(3x^2-4x+2)^2}$ : $f$ est décroissante
    sur $\intFF0{\frac23}$ et croissante sur $\intFF{\frac23}1$, avec $f(0)=\res{\tfrac12}$, $f\left(\tfrac13\right)=\res{0}$,
    $f\left(\tfrac23\right)=\res{-\tfrac12}$ et $f(1)=\res{0}$.
  \item L'angle est droit si $f(x)=0$, soit $\res{x=\tfrac13}$ ou $\res{x=1}$. Il est aigu pour $0\leq x<\frac13$ ($f>0$) et obtus pour
    $\frac13<x<1$ ($f<0$). Comme $MB=ME$, le triangle est équilatéral si et seulement si $\theta=\frac{\pi}3$, c'est-à-dire
    $f(x)=\frac12$ : $3x^2-4x=0$, soit $\res{x=0}$ ($M=D$).
  \item Sur $\intFF0\pi$, la fonction cosinus est strictement décroissante : $\theta(x)$ est d'autant plus grand que $f(x)$ est petit.
    L'angle est maximal pour $x=\frac23$ : $\cos\theta=-\frac12$, $\res{\theta=\tfrac{2\pi}{3}}$ ($120^\circ$). Il est minimal en
    $x=0$, où $f$ atteint son maximum $\frac12$ : $\res{\theta=\tfrac{\pi}{3}}$ ($60^\circ$).
\end{questions}
\end{corrige}

\begin{corrige}{10}{Un angle maximal dans un tétraèdre}
\begin{questions}
  \item $M=B+x\vect{BC}$ avec $\vect{BC}\,(0\pv-1\pv1)$ : $\res{M(0\pv1-x\pv x)}$. Donc $\vect{MA}\,(1\pv x-1\pv-x)$ et
    $\vect{MO}\,(0\pv x-1\pv-x)$.
  \item $\vect{MA}\cdot\vect{MO}=(x-1)^2+x^2=2x^2-2x+1$, $MA^2=2x^2-2x+2$ et $MO^2=2x^2-2x+1$. Ainsi
    $\cos^2\theta(x)=\frac{(2x^2-2x+1)^2}{(2x^2-2x+2)(2x^2-2x+1)}=\frac{2x^2-2x+1}{2x^2-2x+2}$.
  \item $g(0)=g(1)=\frac12$ et $\cos\theta\geq0$ : $\res{\theta(0)=\theta(1)=\tfrac{\pi}{4}}$. Le produit scalaire
    $2x^2-2x+1$ est strictement positif (discriminant $-4<0$), donc $\cos\theta(x)>0$ et $0\leq\theta(x)<\frac\pi2$.
  \item $g(x)=1-\frac{1}{2x^2-2x+2}$ a pour dérivée $\frac{4x-2}{(2x^2-2x+2)^2}$ : $g$ décroît sur $\intFF0{\frac12}$ et croît sur
    $\intFF{\frac12}1$, avec $g\left(\frac12\right)=\frac13$. Sur $\intFF0{\frac\pi2}$, $\theta$ est maximal quand $\cos\theta$ est minimal,
    soit pour $\res{x=\tfrac12}$ : $\cos\theta=\frac{1}{\sqrt3}\approx0{,}577$, d'où $\theta_{\max}\approx\res{54{,}7^\circ}$.
  \item $\vect{MB}\,(0\pv x\pv-x)$ : $\vect{MA}\cdot\vect{MB}=x(x-1)+x^2=x(2x-1)$. Pour $M\neq B$ ($x\neq0$), le triangle est rectangle
    en $M$ si et seulement si $\res{x=\tfrac12}$ : $M$ est le milieu de $[BC]$.
\end{questions}
\end{corrige}

\begin{corrige}{11}{La grande diagonale d'un cube}
\begin{questions}
  \item $\vect{AG}\,(1\pv1\pv1)$, $\vect{BD}\,(-1\pv1\pv0)$ et $\vect{BE}\,(-1\pv0\pv1)$ : $\vect{AG}\cdot\vect{BD}=0$ et
    $\vect{AG}\cdot\vect{BE}=0$, avec $\vect{BD}$ et $\vect{BE}$ non colinéaires. Donc $\vect{AG}$ est normal à $(BDE)$, d'équation
    $x+y+z+d=0$ ; avec $B$, $d=-1$ : $\res{x+y+z-1=0}$.
  \item Les points de $(AG)$ s'écrivent $(t\pv t\pv t)$ : $3t-1=0$, d'où $\res{K\left(\tfrac13\pv\tfrac13\pv\tfrac13\right)}$. C'est
    aussi le point de coordonnées $\left(\frac{1+0+0}{3}\pv\frac{0+1+0}{3}\pv\frac{0+0+1}{3}\right)$ : le centre de gravité de $BDE$.
  \item $(AK)$ est orthogonale au plan : $K$ est le projeté de $A$ et la distance cherchée est $AK=\sqrt{\frac39}=\res{\tfrac{\sqrt3}{3}}$.
  \item Avec la base $ABD$ (aire $\frac12$) et la hauteur $AE=1$ : $V=\frac13\times\frac12\times1=\res{\tfrac16}$. Avec la base $BDE$ et la
    hauteur $AK$ : $\frac16=\frac13\times\mathcal{A}_{BDE}\times\frac{\sqrt3}{3}$, d'où $\mathcal{A}_{BDE}=\res{\tfrac{\sqrt3}{2}}$, qui est
    bien l'aire d'un triangle équilatéral de côté $\sqrt2$ : $\frac{\sqrt3}{4}\times2$.
  \item $C(1\pv1\pv0)$, $F(1\pv0\pv1)$ et $H(0\pv1\pv1)$ vérifient $x+y+z-2=0$ : c'est une équation de $(CFH)$, dont le vecteur normal
    $(1\pv1\pv1)$ est celui de $(BDE)$ : les plans sont parallèles. Ce plan coupe $(AG)$ pour $t=\frac23$. Les points de paramètres
    $0$, $\frac13$, $\frac23$ et $1$ découpent $[AG]$ en $\res{\text{trois segments de longueur }\tfrac{\sqrt3}{3}}$.
\end{questions}
\end{corrige}
